\section{Scope, use rule, and closure state at entry to Phase~IX}\label{sec:p9-scope}

This packet is the project-internal final technical-closure gate required by Phase~IX of the upgrade plan. It does not alter the mathematics already carried by the hardened live proof surface. Instead it assembles, in one place, the final pass conditions, the theorem-node closure matrix, the submission-package inventory, the sign-off checklist, and the closure memo that records why each previously identified blocker is now closed.

\paragraph{Live proof burden.}
The live proof burden remains exactly the four proof-home roots \texttt{core.tex}, \texttt{companion1.tex}, \texttt{companion2.tex}, and \texttt{companion3.tex}, together with the theorem-facing appendices compiled inside those roots. The proof kernel \texttt{proofkernel.tex} and this Phase~IX packet are synchronized audit documents only: they add no mathematics and do not change the proof burden.

\paragraph{Use rule.}
A referee-grade read should proceed in the following order. First attack the core paper or the proof kernel. Second, if a live theorem node needs more detail, descend only to the explicit proof home cited for that theorem. Third, consult this Phase~IX packet only for the final closure matrix, submission inventory, and checklist. Reserve routing notes remain non-mathematical and should never be required to verify a live theorem edge.

\paragraph{Closure state at entry.}
By the time this packet is assembled, the earlier phases have already frozen the theorem surface, repaired the known seam defects, installed theorem-local ledgers, hardened quantifiers and domain transport, exposed the functional / measure backbone, extracted the proof kernel, narrowed the claim framing to the exact local-net / vacuum-sector endpoint, and completed the manuscript-side adversarial-verification pass. The only remaining internal task is therefore to record explicitly that the packaged suite now satisfies the required Phase~IX pass conditions in one synchronized place.
